\documentstyle[% 


righttag% 


,11pt% 


,amssymb% 


,russian%               remove in Eng. 


,izvru%                 Set izven% in Eng. 


% ,izvru-ks             for brief communications 


]{amsart} 


 


%       Math definitions 


 


\newcommand{\re}{\operatorname{Re}} 


\newcommand{\tts}[1]{} 


 


\theoremstyle{plain} 


\theorembodyfont{\it} 


\newtheorem{thmnonum}{\hskip\parindent Теорема} 


\renewcommand{\thethmnonum}{} 


 


\theoremstyle{definition} 


\theorembodyfont{\rm} 


\newtheorem{cornonum}{\hskip\parindent Следствие} 


\newtheorem{notenonum}{\hskip\parindent Замечание} 


\renewcommand{\thecornonum}{} 


\renewcommand{\thenotenonum}{} 


\renewcommand{\thefootnote}{}


%\hoffset=-15mm%                  For Eng. 


 


\setcounter{page}{50} 


\god{2000} 


\nomer{4~(455)} %                        Remove in Eng. 


%\nomer{Vol.\,44, No.\,4}%               Set in Eng. 


%\str{\pageref{begin}--\pageref{end}}%  Set in Eng. 


\udk{517.968} 


 


\author{\it З.Б.\,Цалюк} 


% NB:   ^^ remove in Eng. 


\title{Асимптотическая структура резольвенты неустойчивого уравнения  


Вольтерра\\ с разностным ядром} 


 


\date{} 


% \pagestyle{myheadings}%         Set in Eng. 


\pagestyle{plain} %              Remove in Eng. 


\begin{document} 


% \thispagestyle{plain}%          Set in Eng. 


\maketitle 


\label{begin} 


 


Асимптотическая структура резольвенты $R(t)$ уравнения 


\begin{equation} 


x(t)=\int^t_0 K(t-s)x(s)ds+f(t) 


\end{equation} 


изучалась многими авторами (напр., [1]--[8]). Пусть $K\in L_1$. 


Обозначим через $\wh U(z)=\int\limits^\infty_0 e^{-zt}U(t)dt$ 


преобразование Лапласа 


функции $U$. Так как $\wh R=\frac{\wh K}{1-\wh K}$, то структура $R$ 


существенно 


связана с нулями функции $\wh K(z)-1$ в правой полуплоскости. Если 


$\wh K(z)-1\ne0$ при $\re z\ge0$, то $R\in L_1$ по теореме Винера [1]. 


Если $\wh K(z)-1$ в полуплоскости $\re z\ge0$ имеет конечное число 


нулей $\lambda_r$ целой кратности $m_r$, то, как следует из теории 


вычетов, 


характер поведения $R$ тесно связан с квазимногочленом 


$Q(t)=\sum e^{\lambda_r t}P_{m_r-1}(t)$, где $P_q(t)$ --- многочлен 


степени не выше $q$. 


 


В [2]--[5] функция $R$ представлялась в форме $R=R_0+Q$, $R_0\in L_1$. 


В [6], [7] было предложено представление $R=R_0+Q+Q*R_0$ 


или $R=R_0+Q*R_0$, где $R_0\in L_1$, а 


$$ 


Q*R_0(t)=\int^t_0 Q(t-s)R_0(s)ds. 


$$ 


Отметим, что отсюда получается и представление $R=Q+R_0$. 


 


Так как функция $K(z)$ аналитична при $\re z>0$, то в этой области 


функция 


$\wh K(z)-1$ может иметь нули лишь целой кратности. Однако на мнимой 


оси возможны нули и не целой кратности. В [8] для $K\ge0$ и 


$\wh K(z)-1=z^\beta\psi(z)$, $\psi(0)\ne0$, $\beta\in(0,1)$, была 


получена 


асимптотика $\int\limits^t_0 R(s)ds\sim c t^\beta$. 


 


Данная работа посвящена выяснению структуры резольвенты в 


общем случае конечного числа нулей произвольной кратности. Основным 


условием во всех цитируемых работах является требование $t^p K\in L_1$ 


при некотором 


целом $p$. Если $\wh K(z)-1$ имеет на мнимой оси корень $i\gamma$ 


кратности 


$m>0$, то такое условие представляется естественным, т.\,к. 


в этом случае существует $\wh K^{\,(l)}(i\gamma)=\lim\limits_{z\to 


i\gamma}\wh K^{\,(l)}(z)$, $l$ --- целая 


часть числа $m$, а $\wh K^{\,(l)}(z)$ --- преобразование Лапласа $t^l 


K$. Заметим, 


что если $m<p$, то согласно формуле Тейлора $m$ целое. 


Таким образом, $m$ может быть не целым лишь в случае $t^p K\in L_1$, 


$t^{p+1}K\notin L_1$, 


причем $m=p+\alpha$, $\alpha\in(0,1)$. 


 


Всюду далее через $P_r(t)$ обозначается многочлен степени, не 


превышающей $r$. 


Считаем $P_{-1}(t)=0$ и $\sum\limits^0_1=0$. 


 


\begin{thmnonum} 


Пусть уравнение $\wh K(z)=1$ имеет при $\re z\ge0$ нули 


$\lambda_1=i\gamma_1,\dotsc,\lambda_k=i\gamma_k,\lambda_{k+1},\dotsc, 


\lambda_n$ кратности 


$m_1+\alpha_1,\dotsc,m_k+\alpha_k,m_{k+1},\dotsc,m_n$ 


соответственно, где $m_j\ge0$ целые, 


$\alpha_j\in(0,1)$, причем $m_1=m_2=\dotsb=m_k=p= 


\max\limits_{\re\lambda_j=0}m_j$. 


Пусть далее $t^p K\in L_1$ и 


\begin{equation} 


(t+1)^{p-1+\alpha_j}-(t+1)^{p-1+\alpha_j}* 


e^{-i\gamma_j t}K(t)\in L_1. 


\end{equation} 


Тогда $R(t)=R_1(t)+Q+Q*R_2(t)=R_3+Q*R_4$, где 


$$ 


Q(t)=\sum^k_{j=1}e^{i\gamma_j t}t^{\alpha_j-1}P_p(t)+ 


\sum^n_{j=1}P_{m_j-1}(t)e^{\lambda_j t},\quad R_l\in L_1. 


$$ 


Если, кроме того, $K(t)\to0$ при $t\to\infty$, то и $R_l\to0$ 


при $t\to\infty$. 


\end{thmnonum} 


 


\begin{pf} 


Если $U\in L_1(0,a)$ при любом $a$, то через $\wt U$ 


обозначим действующий в $C[0,\infty)$ оператор свертки, т.\,е. 


оператор, 


определяемый равенством $(\wt U x)(t)=\int\limits^t_0 U(t-s)x(s)ds$, 


а через $I$ --- тождественный оператор. Положим $Q_j(t)=c e^{\lambda_j 


t}$, $j=1,\dotsc,n$, $c>\max\re\lambda_j$, 


$Q_0(t)=c\sum\limits^k_{j=1}t^{\alpha_j-1}e^{i\gamma_j t}$, 


где число $c$ достаточно мало и будет уточнено далее. 


 


Определим ядро $K_0$ равенством $I-\wt K_0=(I-\wt 


K)\prod\limits^n_{j=1}(I+\wt Q_j)^{m_j}$, 


а ядро $K_1$ равенством $I-\wt K_1=(I-\wt K_0)(I+\wt Q_0)$. Покажем, 


что 


$K_1\in L_1(0,\infty)$ и $1-\wh K_1(z)\ne0$ при $\re z\ge0$. 


Справедливы 


соотношения ($r<m_j$) 


\begin{multline*} 


\int^\infty_t e^{\lambda(t-\tau)}\int^\infty_\tau K(s)(\tau-s)^r 


e^{\lambda_j(\tau-s)}ds\,d\tau= \\ 


= 


\begin{cases} 


P_0\displaystyle\int^\infty_t K(s)(t-s)^{r+1} 


e^{\lambda_j(t-s)}ds, &\text{если }\;\lambda=\lambda_j;\\[3mm] 


\displaystyle\int^\infty_t K(s)[P_r(t-s)e^{\lambda_j(t-s)}+ 


P_0e^{\lambda(t-s)}]ds, &\text{если }\;\lambda\ne\lambda_j, 


\end{cases} 


\end{multline*} 


и 


$$ 


\int^\infty_0\int^\infty_t|K(s)|(s-t)^r 


e^{-\re\lambda_j(s-t)}ds\,dt\le\int^\infty_0 


|K(s)|(C_0+C_1 s^p)ds<\infty. 


$$ 


Используя эти соотношения, равенства $\wh K(\lambda_j)=1$, $\wh 


K^{\,(r)}(\lambda_j)=0$, $1\le r<m_j$, 


и индукцию при последовательном умножении на операторы $I+\wt Q_j$, 


легко показать 


$$ 


K_0(t)=K(t)-\sum^n_{j=1}\int^\infty_t K(s)P_{m_j-1} 


(t-s)e^{\lambda_j(t-s)}ds\quad 


\text{и}\quad K_0\in L_1. 


$$ 


По условию 


$$ 


1-\wh K(z)=\prod^k_{j=1}(z-i\gamma_j)^{m_j+\alpha_j} 


\prod_{j>k}(z-\lambda_j)^{m_j}w(z), 


$$ 


где $w(z)\ne0$ при $\re z\ge0$. Так как 


$$ 


1+\wh Q_j(z)=1+\frac{c}{z-\lambda_j}=\frac{z-\lambda_j+c}{z-\lambda_j} 


$$ 


и $z-\lambda_j+c\ne0$ при $\re z\ge0$, то $1-\wh 


K_0(z)=\prod\limits^k_{j=1}(z-i\gamma_j)^{\alpha_j}w_0(z)$, 


где $w_0(z)\ne0$ при $\re z\ge0$. 


Отсюда, в частности, следует 


$$ 


\int^\infty_0e^{-i\gamma_j t}K_0(t)dt= 


\wh K_0(i\gamma_j)=1. 


$$ 


 


Из равенства $I-\wt K_1=(I-\wt K_0)(I+\wt Q_0)$ получаем 


$K_1(t)=K_0(t)-[Q_0(t)-K_0*Q_0(t)]$. 


Поэтому $K_1\in L_1$, если 


$$ 


U_j(t)=\frac{e^{i\gamma_j t}}{t^{1-\alpha_j}}-K_0* 


\frac{e^{i\gamma_j t}}{t^{1-\alpha_j}}\in L_1. 


$$ 


Обозначим $\varphi_j(t)=\int\limits^\infty_t K_0(s)e^{-i\gamma_j s}ds$. 


Так как 


\begin{multline*} 


U_j(t+1)=e^{i\gamma_j(t+1)}\bigg[\frac{1}{(t+1)^{1-\alpha_j}} 


-\int^t_0\frac{K_0(s)e^{-i\gamma_j s}}{(t-s+1)^{1-\alpha_j}}ds


\bigg]+ \\ 


% 


+\int^{t+1}_t\frac{K_0(s)e^{i\gamma_j(t-s+1)}}{(t-s+1)^{1-\alpha_j}} 


ds=e^{i\gamma_j(t+1)}\bigg[\varphi_j(t)- 


\frac{1-\alpha_j}{(t+1)^{2-\alpha_j}}*\varphi_j(t)\bigg] 


+U(t), 


\end{multline*} 


где $U\in L_1$, то $U_j\in L_1$, если 


$\varphi_j-\frac{1-\alpha_j}{(t+1)^{2-\alpha_j}}*\varphi_j\in L_1$. 


Но 


\begin{multline*} 


\varphi_j(t)=c\int^\infty_t e^{-i\gamma_j s}K(s)ds- 


\sum_{l\ne j}\int^\infty_t e^{-i\gamma_j s}K(s) 


P_{m_l-1}(t-s)e^{(\lambda_l-i\gamma_j)(t-s)}ds - \\ 


% 


-\int^\infty_t e^{-i\gamma_j s}K(s)P_p(t-s)ds. 


\end{multline*} 


Как было отмечено выше, $\sum\limits_{l\ne j}\in L_1$ и, если $p>0$, 


$$ 


\int^\infty_t e^{-i\gamma_j s}K(s)P_{p-1}(t-s)ds\in L_1\quad 


\text{и}\quad 


\int^\infty_t e^{i\gamma_j s}K(s)ds\in L_1. 


$$ 


Таким образом, 


$$ 


\varphi_j(t)=c\int^\infty_t(s-t)^p e^{-i\gamma_j s}K(s)ds+ 


U(t)=c\psi_j(t)+U(t), 


$$ 


где $U\in L_1$. Так как $\wh K(i\gamma_j)=1$, $\wh 


K^{\,(l)}(i\gamma_j)=0$, $l=1,\dotsc,p$, 


то $\psi(0)=\dotsb=\psi^{(p-1)}(0)=0$, $\psi^{(p)}(0)=p!$ Используя эти 


равенства, интегрированием по частям найдем 


$$ 


\psi_j(t)-\frac{1-\alpha_j}{(t+1)^{2-\alpha_j}}*\psi_j=c 


[(t+1)^{p-1+\alpha_j}-(t+1)^{p-1+\alpha_j} 


*e^{-i\gamma_j t}K(t)]+\sum^{p-1}_{k=0}c_k 


\int^\infty_t(s-t)^k e^{-i\gamma_j s}K(s)ds. 


$$ 


По условию $t^p K\in L_1$. Поэтому $\int\limits^\infty_t(s-t)^k 


e^{-i\gamma_j s}K(s)ds\in L_1$. 


Следовательно, $\psi_j-\frac{1-\alpha_j}{(t+1)^{2-\alpha_j}}*\psi_j\in 


L_1$ и 


$$ 


\varphi_j-\frac{1-\alpha_j}{(t+1)^{2-\alpha_j}}*\varphi_j= 


c\bigg(\psi_j-\frac{1-\alpha_j}{(t+1)^{2-\alpha_j}}*\psi_j\bigg)+ 


\bigg(U-\frac{1-\alpha_j}{(t+1)^{2-\alpha_j}}*U\bigg)\in L_1. 


$$ 


Итак, $K_1\in L_1$. Покажем, что $1-\wh K_1(z)\ne0$ при $\re z\ge0$. 


Имеем 


$$ 


1-\wh K_1(z)=(1-\wh K_0(z))(1+\wh Q_0(z))= 


\prod^k_{j=1}(z-i\gamma_j)^{\alpha_j}w_0(z)\bigg[ 


1+c\sum^k_{j=1}\frac{\Gamma(\alpha_j)}{(z-i\gamma_j)^{\alpha_j}} 


\bigg]=w_0(z)w_1(t). 


$$ 


Покажем также, что $w_1(z)\ne0$ в полуплоскости $\re z\ge0$. Обозначим 


$B_j(c)=\{z:|z-i\gamma_j|\le c,\ \re z\ge0\}$. 


Если $z\notin\cup B_j(c)$, то 


$$ 


|w_1(z)|\ge\prod_j|z-i\gamma_j|^{\alpha_j}\bigg[1-c\sum_j 


\frac{\Gamma(\alpha_j)}{|z-i\gamma_j|^{\alpha_j}}\bigg]\ge 


c^{\alpha_1+\dotsb+\alpha_k}\bigg[1-\sum_j 


\Gamma(\alpha_j)c^{1-\alpha_j}\bigg]>0, 


$$ 


если $c$ достаточно мало. Пусть $z\in B_1(c)$. Обозначим 


$$ 


f(z)=\prod^k_{j=1}(z-i\gamma_j)^{\alpha_j}+c\Gamma(\alpha_1) 


\prod^k_{j=2}(z-i\gamma_j)^{\alpha_j},\quad 


g(z)=c\sum^k_{j=2}\Gamma(\alpha_j)\prod_{l\ne j} 


(z-i\gamma_l)^{\alpha_l}. 


$$ 


Так как $|\arg(z-i\gamma_1)|\le\pi/2$, то 


$|\arg(z-i\gamma_1)^{\alpha_1}|<\pi/2$. Следовательно, 


$\re(z-i\gamma_1)^{\alpha_1}\ge 0$ и 


$|(z-i\gamma_1)^{\alpha_1}+c\Gamma(\alpha_1)|\ge c\Gamma(\alpha_1)$. 


Поэтому при $c<\frac{1}{2}\min\limits_{j\ne l}|\gamma_j-\gamma_l|=d$ 


имеем 


\begin{align*} 


|f(z)|&=\prod^k_{j=2}|z-i\gamma_j|^{\alpha_j}|(z-i\gamma_1)^{\alpha_1}+ 


c\Gamma(\alpha_1)|\ge d^{\alpha_2+\dotsb+\alpha_k}\Gamma(\alpha_1)c,\\ 


% 


|g(z)|&=c|z-i\gamma_1|^{\alpha_1}\bigg|\sum^k_{j=2} 


\Gamma(\alpha_j)\prod_{l\ne j,1}(z-i\gamma_l)^{\alpha_l}\bigg|\le 


M c^{1+\alpha_1}. 


\end{align*} 


Следовательно, для $z\in B_1(c)$ при достаточно малом $c$ 


$$ 


|w_1(z)|\ge|f(z)|-|g(z)|\ge 


d^{\alpha_2+\dotsb+\alpha_k}\Gamma(\alpha_1)c- 


M c^{1+\alpha_1}>0. 


$$ 


Аналогично доказывается, что $|w_1(z)|>0$ при $z\in B_j(c)$. Таким 


образом, $w_1(z)\ne0$ при $\re z\ge0$, а значит, и $1-\wh K_1(z)\ne0$ 


при $\re z\ge0$. 


 


Положим $I+\wt D_0=\prod\limits^n_{j=1}(I+\wt Q_j)^{m_j}$ и $I+\wt 


D_1=(I+\wt D_0)(I+\wt Q_0)$. 


Ядро оператора $\wt D_0$ есть сумма сверток всевозможных ядер $Q_j$, и 


легко 


подсчитать, что $D_0(t)=\sum\limits^n_{j=1}P_{m_j-1}(t)e^{\lambda_j 


t}$. 


Интегрированием по частям получим 


$$ 


\int^t_0\frac{e^{\lambda(t-s)}(t-s)^r e^{i\gamma s}}{s^{1-\alpha}} 


ds= 


\begin{cases} 


e^{\lambda t}P_r(t)+c e^{i\gamma t}t^{\alpha-1}+U(t), 


\text{ где }U\in L_1\cap C_0, &\text{если } 


\lambda\ne i\gamma;\\ 


c e^{i\gamma t}t^{r+\alpha}, &\text{если }\lambda=i\gamma,


\end{cases} 


$$ 


где $C_0$ --- пространство непрерывных функций $x$, имеющих


$\lim\limits_{t\to\infty}x(t)=0$.


Следовательно,


$$ 


D_1(t)=D_0+Q_0+D_0*Q_0(t)=\sum^k_{j=1}e^{i\gamma_j t} 


t^{\alpha_j-1}P_p(t)+\sum^n_{j=1}P_{m_j-1}(t) 


e^{\lambda_j t}+v(t)=Q+v, 


$$ 


где $v\in L_1$ и $v\to0$ при $t\to\infty$. Пусть $R_1$ --- резольвента 


ядра $K_1$. 


Так как $K_1\in L_1$ и $1-\wh K_1(z)\ne0$ при $\re z\ge0$, то по 


теореме Винера [1]\; $R_1\in L_1$. Из $I-\wt K_1=(I-\wt K)(I+\wt D_1)$ 


следует 


$I+\wt R=(I+\wt D_1)(I+R_1)$ и, значит,


$R=D_1+R_1+D_1*R_1=R_1+v+v*R_1+Q+Q*R_1=R_2+Q+Q*R_1$, 


где $R_2=R_1+v+v*R_1\in L_1$, 


т.\,к. $v,R_1\in L_1$. Кроме того, если $K\in C_0$, то $K_0\in C_0$, а 


т.\,к. 


$Q_0\in C_0$, $K_0\in L_1$, то и $K_1=K_0-Q_0+K_0*Q_0\in C_0$. Но 


$R_1=K_1+R_1*K_1$ и $R_1\in L_1$, поэтому $R_1*K_1\in C_0$ и $R_1\in 


C_0$. 


Следовательно, и $R_2\in C_0$. 


 


Второе представление $R$ получим следующим образом. Определим 


ядро $R_0(t)$ равенством $\wt R=(I+\wt D_1)\wt R_0$. Из $\wt K=(I-\wt 


K)\wt R=(I-\wt K_1)\wt R_0$ 


следует $\wt R_0=(I+\wt R_1)\wt K$ и, значит, $R_0\in L_1$. 


Таким образом, $R(t)=R_0(t)+D_1*R_0(t)=R_0(t)+v*R_0(t)+Q*R_0(t)$. 


Кроме того, если $K\in C_0$, то и $R_0\in C_0$, а т.\,к. $R_0\in L_1$ и 


$v\in C_0$, то $R_0+v*R_0\in C_0$. 


\end{pf} 


 


\begin{notenonum} 


Как следует из доказательства, условие (2) может быть 


заменено на условие 


$$ 


\psi_j(t)-\frac{1-\alpha_j}{(t+1)^{2-\alpha_j}}*\psi_j\in L_1,


$$ 


где $\psi_j(t)=\int\limits^\infty_t(s-t)^p e^{-i\gamma_j s}K(s)ds$. 


\end{notenonum} 


 


Если $k=0$, т.\,е. все корни целой кратности, то условие (2) 


отсутствует. 


Поэтому возникает вопрос: насколько оно необходимо, если $k\ne0$. 


Для простоты рассмотрим случай $k=n=1$, $m=0$. При действительном 


ядре $K\in L_1$ это означает, что $\wh K(z)-1=z^\alpha\psi(z)$, где 


$\psi(z)\ne0$ при $\re z\ge0$. Тогда условие (2) принимает вид 


\begin{equation} 


\frac{1}{(t+1)^{1-\alpha}}- 


\frac{1}{(t+1)^{1-\alpha}}*K\in L_1, 


\end{equation} 


а резольвента представима в форме 


\begin{equation} 


R=R_1+\frac{c}{t^{1-\alpha}}+\frac{c}{t^{1-\alpha}}*R_1,\quad 


\text{где}\quad 


R_1\in L_1. 


\end{equation} 


Покажем, что если справедливо представление (4), то выполнено и условие 


(3). Действительно, легко проверить, что если выполнено (4), то $R_1$ 


--- резольвента ядра 


$K_1=K-\frac{c}{t^{1-\alpha}}+K*\frac{c}{t^{1-\alpha}}$. Из равенства 


$R_1=K_1+K_1*R_1$ следует $(1+\wh R_1)(1-\wh K_1)=1$. Так как 


$$ 


1-\wh K_1=(1-\wh K)(1+\wh{c t}^{\,\alpha-1})= 


-z^\alpha\psi(z)\frac{z^\alpha+c\Gamma(\alpha)}{z^\alpha} 


\ne\infty 


$$ 


при $\re z\ge0$, то $1+\wh R_1(z)\ne0$ 


в правой полуплоскости. По теореме 


Пэли--Винера [1]\; $-K_1\in L_1$ как резольвента ядра $-R_1$. По 


условию $K\in L_1$, 


следовательно, 


$U(t)=\frac{1}{t^{1-\alpha}}-K*\frac{1}{t^{1-\alpha}}\in L_1$. Осталось 


заметить, что 


$$ 


\frac{1}{(t+1)^{1-\alpha}}-K*\frac{1}{(t+1)^{1-\alpha}}= 


U(t+1)-\int^{t+1}_t\frac{K(s)}{(t-s+1)^{1-\alpha}}ds 


$$ 


и 


$$ 


\int^{t+1}_t\frac{K(s)}{(t-s+1)^{1-\alpha}}ds\in L_1. 


$$ 


 


\begin{cornonum} 


Пусть $\wh K(z)-1=z^{p+\alpha}\psi(z)$, $p\ge0$ целое, 


$\alpha\in(0,1)$, $\psi(z)\ne0$ при $\re z\ge0$. Пусть, кроме того,


$t^p K\in L_1$ и 


\begin{equation} 


(t+1)^{p-1+\alpha}-(t+1)^{p-1+\alpha}*K\in L_1. 


\end{equation} 


Если $f(t)=f_\infty+o(1)$, $t\to\infty$, и $x(t)$ --- соответствующее 


ему решение (1), то 


\begin{equation} 


x(t)=t^{p+\alpha}(c+o(1)),\quad t\to\infty. 


\end{equation} 


Утверждение вытекает из равенств 


\begin{multline*} 


x(t)=f(t)+R*f(t)=f(t)+R_3*f(t)+P_p(t)t^{\alpha-1} 


*(R_4*f(t))= \\ 


% 


= C_1+o(1)+P_p(t)t^{\alpha-1}*(C_2+o(1))=c t^{p+\alpha} 


+o(t^{p+\alpha}). 


\end{multline*} 


Теорема Феллера [8] гласит, что 


$$ 


\int^t_0 R(s)ds=t^\alpha(C_1+o(1)), 


$$ 


если $K\ge 0$, 


$$ 


\int^\infty_0 K(t)dt=1,\quad 


\int^t_0 s K(s)ds=t^{1-\alpha}(C+o(1)). 


$$ 


Из условий теоремы Феллера легко следует, что $\wh 


K(z)-1=z^\alpha\psi(z)$, 


так что при дополнительном условии (5) (с $p=0$) утверждение теоремы 


вытекает из (6). 


 


Не ясно, является ли условие (5) с $p=0$ следствием условий 


теоремы Феллера или в условиях теоремы Феллера для некоторых $R$ 


отсутствует представление (4). 


\end{cornonum} 


 


\begin{thebibliography}{99} 


 


\bibitem{1} 


Винер Н., Пэли Р. {\em Преобразование Фурье в комплексной области}. 


-- М.: Наука, 1964. -- 267 с. 


\bibitem{2} 


Smith W.L. {\em Asymptotic renewel theorems} // Proc.


Roy. Soc. Edinburgh. Sect.\,A. -- 1954. -- V.\,64. -- P.\,9--48. 


\bibitem{3} 


Севастьянов Б.А. {\em Ветвящиеся процессы}. -- М.: Наука, 1971. -- 436 


с. 


\bibitem{4} 


Дербенев В.А. {\em Асимптотика решения неустойчивого уравнения 


восстановления}. -- Ред. журн. ``Изв. вузов. Математика''. -- Казань,


1976. -- 13 с. -- Деп. в ВИНИТИ 01.07.76, \No\,2460--76. 


\bibitem{5} 


Miller R.K. {\em Structure of solutions of unstable linear 


Volterra integrodifferential equations} // J. Different. Equat. -- 


V.\,15. -- P.\,129--157. 


\bibitem{6} 


Дербенев В.А., Цалюк З.Б. {\em Асимптотика резольвенты неустойчивого 


уравнения Вольтерра с разностным ядром} // Матем. заметки. -- 1997. -- 


Т.\,62. -- Вып.\,1. -- С.\,88--94. 


\bibitem{7} 


Цалюк З.Б. {\em Асимптотическое представление резольвенты системы 


интегро-диф\-фе\-рен\-ци\-альных уравнений} // Понтрягинские чтения. Тез. докл. 


-- 


Воронеж, 1996. -- С.\,185. 


\bibitem{8} 


Феллер В. {\em Введение в теорию вероятностей и ее приложения}. 


T.\,2. -- 


М.: Мир, 1967. -- 751~с. 


 


\end{thebibliography} 


 


\vskip 2cm 


 


\post{\it Кубанский государственный университет}{\it Поступила} 


\post{}{10.03.1998} 


 


\label{end} 


\end{document} 


 


